\documentclass[uplatex,a4paper,11pt]{jsarticle}
\usepackage[margin=22mm]{geometry}
\usepackage[dvipdfmx]{hyperref,xcolor}
\usepackage{fvextra,enumitem,longtable,booktabs,amssymb}
\hypersetup{
  colorlinks=true,
  linkcolor=blue!55!black,
  urlcolor=blue!65!black,
  pdftitle={TeXNote P0リバースプロキシ・VLAN・ファイアウォール構築手順}
}
\definecolor{commandbg}{RGB}{244,246,248}
\DefineVerbatimEnvironment{command}{Verbatim}{
  fontsize=\small,
  frame=single,
  rulecolor=\color{gray!55},
  breaklines=true,
  breakanywhere=true
}
\setlist{nosep}

\title{TeXNote P0リバースプロキシ・VLAN・\
ファイアウォール構築手順\
\large P1上での先行構築からP0分離まで}
\author{}
\date{2026年8月24日}

\begin{document}
\maketitle

\begin{abstract}
本書は、TeXNoteのP1認証サーバーへNginxリバースプロキシとUFWを先行導入し、
将来VLAN対応ルーターとマネージドスイッチを購入した後、公開入口を専用Raspberry Pi
「P0」へ移す手順をまとめたものである。P1はユーザー認証とP2中継、P2はTeXコンパイル、
P0はHTTPS終端と外部公開境界を担当する。作業は段階ごとに動作確認し、問題があれば
直前の構成へ戻せるように進める。
\end{abstract}

\tableofcontents
\newpage

\section{目的と最終構成}

\subsection{現在から最終形への移行}

最初はP1上でNginxを動かし、リバースプロキシを経由するアプリ構成を先に完成させる。

\begin{command}
段階1（LAN内試験）

TeXNote
   | HTTP :80
   v
P1 Nginx
   | HTTP 127.0.0.1:8000
   v
P1 FastAPI
   | 内部トークン付きHTTP
   v
P2 FastAPI
\end{command}

VLAN対応機器とP0を準備した後、NginxとHTTPS証明書をP0へ移す。

\begin{command}
最終構成

Internet
   | HTTPS :443
   v
P0 Nginx（DMZ VLAN）
   | HTTP :8000（P0からだけ許可）
   v
P1 FastAPI（認証 VLAN）
   | HTTP :8000（P1からだけ許可）
   v
P2 FastAPI（コンパイル VLAN）
\end{command}

外部へ公開するのはP0のTCP 443だけである。P1、P2、SSH、TCP 8000をインターネットへ
直接公開しない。

\subsection{各機器の役割}

\begin{longtable}{@{}p{18mm}p{46mm}p{82mm}@{}}
\toprule
機器 & 役割 & 保持するもの・実行するもの\\
\midrule
P0 & 公開境界 & Nginx、HTTPS証明書、32 MiB制限、アクセスログ\\
P1 & 認証・API中継 & FastAPI、JWT、SQLiteユーザーDB、利用制限、P2内部トークン\\
P2 & TeXコンパイル & FastAPI、TeX Live、添付検査、隔離した版組処理\\
\bottomrule
\end{longtable}

\section{作業前の原則}

\begin{enumerate}
  \item 現行サービスの状態と設定ファイルを記録する。
  \item 新しい経路が動くことを確認してから古い経路を閉じる。
  \item SSH許可規則を追加してからUFWを有効にする。
  \item P1上のNginxが完成するまでTCP 8000の待受方法を変更しない。
  \item P0移行が完了するまでP1上のNginx設定を削除しない。
  \item VLAN間は既定拒否とし、必要な送信元、宛先、ポートだけを許可する。
  \item IPv4だけでなくIPv6側の公開状態も確認する。
\end{enumerate}

以下の現行値を前提とする。異なる場合は読み替える。

\begin{longtable}{@{}p{52mm}p{96mm}@{}}
\toprule
項目 & 現行値\\
\midrule
P1 LANアドレス & \texttt{192.168.3.21}\\
P2 LANアドレス & \texttt{192.168.3.4}\\
P1 systemdサービス & \texttt{tex-auth.service}\\
P1プログラム & \texttt{/home/mat/tex-auth}\\
P1仮想環境 & \texttt{/home/mat/tex-auth/.auth-env}\\
P1 FastAPIポート & TCP 8000\\
Nginx LAN試験ポート & TCP 80\\
公開時HTTPSポート & TCP 443\\
HTTP要求上限 & P0、P1、P2すべて32 MiB\\
\bottomrule
\end{longtable}

\newpage
\section{全体の実施順序}

本書とP1構築手順書に分散している作業を、実機へ適用する順番にまとめる。
各段階では、確認項目がすべて成功してから次へ進む。失敗した場合は次の段階へ進まず、
その段階で退避した設定へ戻す。

\subsection{1：P1ソースとプラン管理コマンドを更新する}

最初に、P1の現行プログラムとSQLite DBを退避する。P1上で実行する。

\begin{command}
cd /home/mat/tex-auth
backup_dir=/home/mat/tex-auth-backup-before-plan
mkdir -p "$backup_dir"
cp -p main.py auth.py database.py tex_service.db "$backup_dir"/
test ! -f plans.py || cp -p plans.py "$backup_dir"/
test ! -d tools || cp -a tools "$backup_dir"/
\end{command}

管理用MacのTeXNoteリポジトリ直下から、秘密値とDBを含めずに現行ソースを転送する。

\begin{command}
scp Server/P1/progs/main.py \
    Server/P1/progs/auth.py \
    Server/P1/progs/database.py \
    Server/P1/progs/plans.py \
    mat@192.168.3.21:/home/mat/tex-auth/

scp -r Server/P1/progs/tools \
    mat@192.168.3.21:/home/mat/tex-auth/
\end{command}

P1上で構文を検査してからサービスを再起動する。既存DBに対して
\texttt{tools.init\_db}を実行し直してはならない。

\begin{command}
cd /home/mat/tex-auth
source .auth-env/bin/activate
python -m py_compile *.py tools/*.py
sudo systemctl restart tex-auth.service
systemctl status tex-auth.service --no-pager -l
journalctl -u tex-auth.service -n 100 --no-pager
\end{command}

\subsection{2：管理者コマンドでplan変更を確認する}

既存ユーザーを例としてStandardへ変更し、DBの結果を確認する。メールアドレスは実際の
登録ユーザーへ置き換える。

\begin{command}
cd /home/mat/tex-auth
source .auth-env/bin/activate
python -m tools.change_plan user@example.com standard
python -m tools.check02 user@example.com
\end{command}

\texttt{change\_plan}が変更前と変更後のplanを表示し、\texttt{check02}にも
\texttt{standard}が表示されることを確認する。変更後のplanは要求ごとにDBから読み込むため、
plan変更だけを理由にP1を再起動する必要はない。現段階では\texttt{/me}を利用した
プラン変更UIは作らず、管理者コマンドだけを正式な変更経路とする。

\subsection{3：P1上のNginx経由を完成させる}

本書の「段階1：P1上にNginxを構築する」に従う。Uvicornの待受を変更する前に、
Nginx経由の\texttt{/login}、\texttt{/me}、\texttt{/compile}を確認する。
同時に、Nginx経由の\texttt{/compile-test}、\texttt{/docs}、
\texttt{/openapi.json}が404になることを確認する。

Uvicornが\texttt{127.0.0.1:8000}だけで待ち受けていることを確認する。実機が既に
この状態なら設定編集、\texttt{daemon-reload}、待受変更を目的としたP1再起動は行わない。
\texttt{0.0.0.0:8000}または\texttt{[::]:8000}で待ち受けている場合だけ、Nginx経由の
正常動作を確認した後に\texttt{127.0.0.1:8000}限定へ変更する。
\texttt{/compile-test}は削除せず、P1へSSH接続した管理者が
\texttt{127.0.0.1:8000/compile-test}を使う保守APIとして残す。

\subsection{4：P1のUFWを適用する}

SSHとNginxへの許可規則を先に追加してからUFWを有効にする。別ターミナルから新しい
SSH接続が成功するまで、作業中のSSH接続を閉じない。LAN上の端末からP1のTCP 8000へ
直接接続できず、P1自身から\texttt{127.0.0.1:8000}へ接続できることを確認する。
その後P1を再起動し、Nginx、P1、UFWが自動復旧することを確認する。

\subsection{5：P0とVLANへ公開境界を移す}

段階1が安定してから、本書の「段階2：P0を準備する」以降へ進む。P0のNginxは
\texttt{/login}、\texttt{/me}、\texttt{/compile}だけをP1へ転送する。
P1のUvicornは認証VLAN側アドレスだけで待ち受け、P1のUFWはP0からのTCP 8000だけを
許可する。P2のUFWはP1からのTCP 8000だけを許可する。

P0経由の通常版組まで成功し、P0からP2へ直接接続できないことを確認するまで、P1上の
Nginx設定を削除しない。問題が発生した場合は、本書の切戻し手順で段階1へ戻す。

\subsection{6：HTTPSを設定して外部公開する}

最後にP0へHTTPS証明書を設定する。ルーターでインターネットから転送するのは
P0のTCP 443だけとし、TCP 22、80、8000、P1、P2を直接公開しない。
外部から\texttt{/compile-test}、\texttt{/docs}、\texttt{/openapi.json}が
404となり、アプリの\texttt{/login}、\texttt{/me}、\texttt{/compile}だけが
HTTPSで利用できることを最終確認する。

\newpage
\section{段階1：P1上にNginxを構築する}

\subsection{現状確認とバックアップ}

P1上で実行する。

\begin{command}
systemctl status tex-auth.service --no-pager -l
systemctl is-enabled tex-auth.service
systemctl is-active tex-auth.service
sudo ss -ltnp | grep -E ':8000|:80|:443'
\end{command}

現在のサービスファイルを退避する。

\begin{command}
sudo cp -p \
  /etc/systemd/system/tex-auth.service \
  /etc/systemd/system/tex-auth.service.before-nginx
\end{command}

\subsection{Nginxの導入}

\begin{command}
sudo apt update
sudo apt install nginx
nginx -v
sudo systemctl status nginx --no-pager
\end{command}

初期ページを確認する。

\begin{command}
curl -I http://192.168.3.21/
\end{command}

\subsection{P1用Nginx設定}

\begin{command}
sudoedit /etc/nginx/sites-available/texnote-p1
\end{command}

次の内容を設定する。\texttt{proxy\_pass}の末尾にはスラッシュを付けず、
\texttt{/login}、\texttt{/me}、\texttt{/compile}の完全一致だけをP1へ渡す。
\texttt{location /}でその他を404にするため、保守用の\texttt{/compile-test}、
\texttt{/docs}、\texttt{/openapi.json}はNginx経由で公開されない。

\begin{command}
server {
    listen 80;
    listen [::]:80;

    server_name _;
    server_tokens off;

    client_max_body_size 32m;
    client_body_timeout 30s;

    access_log /var/log/nginx/texnote-p1-access.log;
    error_log  /var/log/nginx/texnote-p1-error.log;

    location = /login {
        proxy_pass http://127.0.0.1:8000;
        include /etc/nginx/proxy_params;
        proxy_connect_timeout 5s;
        proxy_send_timeout 30s;
        proxy_read_timeout 30s;
    }

    location = /me {
        proxy_pass http://127.0.0.1:8000;
        include /etc/nginx/proxy_params;
        proxy_connect_timeout 5s;
        proxy_send_timeout 30s;
        proxy_read_timeout 30s;
    }

    location = /compile {
        proxy_pass http://127.0.0.1:8000;
        include /etc/nginx/proxy_params;
        proxy_connect_timeout 5s;
        proxy_send_timeout 120s;
        proxy_read_timeout 120s;
    }

    location / {
        return 404;
    }
}
\end{command}

P0、P1、P2の要求上限を32 MiBへ統一する。Nginxの既定値は1 MiBであるため、
\texttt{client\_max\_body\_size 32m;}を省略しない。

\subsection{サイトの有効化}

\begin{command}
sudo ln -s \
  /etc/nginx/sites-available/texnote-p1 \
  /etc/nginx/sites-enabled/texnote-p1
sudo unlink /etc/nginx/sites-enabled/default
sudo nginx -t
sudo systemctl reload nginx
sudo systemctl enable nginx
\end{command}

\texttt{nginx -t}が\texttt{syntax is ok}、\texttt{test is successful}を返してから
reloadする。

\subsection{Nginx経由とUvicorn待受の確認}

最初に、現在のsystemd設定と実際の待受アドレスを確認する。

\begin{command}
systemctl cat tex-auth.service
sudo ss -ltnp | grep ':8000'
\end{command}

実機が既に\texttt{127.0.0.1:8000}だけで待ち受けている場合は、それが正しい現行状態で
ある。待受設定を変更せず、その状態のままNginx経由を確認する。

\begin{command}
curl --fail-with-body http://127.0.0.1:8000/openapi.json >/dev/null
test "$(curl -sS -o /dev/null -w '%{http_code}' \
  http://192.168.3.21/openapi.json)" = 404
test "$(curl -sS -o /dev/null -w '%{http_code}' -X POST \
  http://192.168.3.21/compile-test)" = 404
\end{command}

最初の要求はP1自身の保守APIを直接確認する。後の2要求はNginxが公開対象外の
パスをP1へ転送せず、認証処理より前に404を返すことを確認する。

登録済みユーザーでログインする。

\begin{command}
curl -sS -X POST http://192.168.3.21/login \
  -H 'Content-Type: application/json' \
  -d '{
    "email":"<登録済みメールアドレス>",
    "password":"<パスワード>"
  }'
\end{command}

\subsection{必要な場合だけUvicornをlocalhost限定に変更}

前項の確認で\texttt{127.0.0.1:8000}だけが表示された場合、この項の編集操作はすべて
省略する。\texttt{0.0.0.0:8000}または\texttt{[::]:8000}が表示された場合に限り、
Nginx経由のログインが成功してから変更する。

\begin{command}
sudoedit /etc/systemd/system/tex-auth.service
\end{command}

\texttt{ExecStart}を1行で次のようにする。

\begin{command}
ExecStart=/home/mat/tex-auth/.auth-env/bin/uvicorn main:app --host 127.0.0.1 --port 8000
\end{command}

設定を編集した場合だけ反映する。

\begin{command}
sudo systemctl daemon-reload
sudo systemctl restart tex-auth.service
systemctl status tex-auth.service --no-pager -l
journalctl -u tex-auth.service -n 100 --no-pager
\end{command}

編集を省略した場合も含め、最後に待受状態を確認する。

\begin{command}
sudo ss -ltnp | grep -E ':8000|:80|:443'
\end{command}

期待する待受状態は次のとおりである。

\begin{command}
0.0.0.0:80       nginx
[::]:80          nginx
127.0.0.1:8000   uvicorn
\end{command}

\texttt{0.0.0.0:8000}が残っていてはいけない。

\subsection{TeXNoteの接続先変更}

macOS版、iPad版、iPhone版のP1 URLを次のように変更する。

\begin{command}
変更前: http://192.168.3.21:8000
変更後: http://192.168.3.21
\end{command}

80番はHTTPの標準ポートなので省略できる。各アプリでログイン、\texttt{/me}、通常版組、
画像付き版組、TeXエラー表示を確認する。

\subsection{ログ確認}

\begin{command}
sudo tail -n 100 /var/log/nginx/texnote-p1-access.log
sudo tail -n 100 /var/log/nginx/texnote-p1-error.log
journalctl -u nginx -n 100 --no-pager
\end{command}

リアルタイム監視は次のとおり。

\begin{command}
sudo tail -f \
  /var/log/nginx/texnote-p1-access.log \
  /var/log/nginx/texnote-p1-error.log
\end{command}

終了は\texttt{Ctrl+C}で行う。

\newpage
\section{段階1：P1のUFWを設定する}

\subsection{UFWの導入}

\begin{command}
sudo apt update
sudo apt install ufw
ufw --version
sudo systemctl status ufw --no-pager
\end{command}

\subsection{有効化前にSSHとNginxを許可}

LANが\texttt{192.168.3.0/24}の場合の例を示す。現在のSSH接続を閉じずに作業する。

\begin{command}
sudo ufw allow from 192.168.3.0/24 to any port 22 proto tcp
sudo ufw allow from 192.168.3.0/24 to any port 80 proto tcp
sudo ufw default deny incoming
sudo ufw default allow outgoing
sudo ufw deny 8000/tcp
sudo ufw show added
\end{command}

内容を確認してから有効化する。

\begin{command}
sudo ufw enable
sudo ufw status verbose
sudo ufw status numbered
\end{command}

Macの別ターミナルから新しいSSH接続を試す。

\begin{command}
ssh mat@192.168.3.21
\end{command}

新しいSSH接続に成功するまで、元のSSH接続を終了しない。

\subsection{UFW適用後の確認}

Nginx経由の\texttt{/login}がP1へ到達する必要がある。空JSONに対する422は、
NginxではなくP1の入力検証が応答したことを示す。

\begin{command}
test "$(curl -sS -o /dev/null -w '%{http_code}' -X POST \
  -H 'Content-Type: application/json' -d '{}' \
  http://192.168.3.21/login)" = 422
test "$(curl -sS -o /dev/null -w '%{http_code}' \
  http://192.168.3.21/openapi.json)" = 404
\end{command}

LANからUvicornへ直接接続する次のコマンドは失敗するのが正常である。

\begin{command}
curl --connect-timeout 5 http://192.168.3.21:8000/docs
\end{command}

P1自身からlocalhostのUvicornへは接続できる。

\begin{command}
curl --fail-with-body \
  http://127.0.0.1:8000/openapi.json >/dev/null
\end{command}

\subsection{再起動試験}

\begin{command}
sudo reboot
\end{command}

再接続後に確認する。

\begin{command}
systemctl is-active nginx
systemctl is-enabled nginx
systemctl is-active tex-auth.service
systemctl is-enabled tex-auth.service
sudo ufw status verbose
sudo ss -ltnp | grep -E ':8000|:80|:443'
curl --fail-with-body http://127.0.0.1:8000/openapi.json >/dev/null
\end{command}

\section{段階1の切戻し}

Nginx経由で問題が発生した場合、バックアップしたsystemd設定を戻す。

\begin{command}
sudo cp -p \
  /etc/systemd/system/tex-auth.service.before-nginx \
  /etc/systemd/system/tex-auth.service
sudo systemctl daemon-reload
sudo systemctl restart tex-auth.service
sudo systemctl disable --now nginx
\end{command}

UFWの8000拒否規則を番号で確認して削除する。

\begin{command}
sudo ufw status numbered
sudo ufw delete <8000/tcp拒否規則の番号>
\end{command}

TeXNoteのURLを従来値へ戻す。

\begin{command}
http://192.168.3.21:8000
\end{command}

\newpage
\section{VLAN機器の安価な購入案}

\subsection{推奨する最小構成}

2026年8月時点で日本国内から入手しやすい構成例を示す。価格と販売状況は購入時に
再確認する。

\begin{longtable}{@{}p{40mm}p{55mm}p{50mm}@{}}
\toprule
機器 & 候補 & 価格目安\\
\midrule
VLANルーター & TP-Link ER605 V2 & 約9,500～11,500円\\
5ポートVLANスイッチ & TP-Link TL-SG105E & 約3,000～3,500円\\
8ポートVLANスイッチ & TP-Link TL-SG108E & 約5,000～6,500円\\
\bottomrule
\end{longtable}

ER605は802.1Q VLAN、複数ネットワーク、送信元・宛先IPによるACL、ファイアウォールに
対応する。TL-SG105E/TL-SG108Eは802.1QタグVLANに対応する。

P0、P1、P2、管理用Mac、ルーターとのtrunkで5ポートをすべて使用するため、長期運用では
8ポート版が扱いやすい。最小費用を優先する場合は5ポート版でも構築できる。

\subsection{購入前に既存ルーターを確認}

既存ルーターが次のすべてへ対応する場合、ER605を追加せずVLANスイッチだけで構築できる。

\begin{itemize}
  \item 802.1Q VLAN
  \item VLANごとのIPインターフェースとDHCP
  \item VLAN間ACL
  \item 送信元、宛先、プロトコル、ポートによる許可・拒否
  \item ポート転送
\end{itemize}

家庭用ルーターの「VLAN」がIPTV専用機能だけの場合は使用できない。

\subsection{参考資料}

\begin{sloppypar}
\begin{itemize}
  \item ER605：\url{https://www.tp-link.com/jp/business-networking/vpn-router/er605/v2/}
  \item TL-SG105E：\url{https://www.tp-link.com/jp/business-networking/easy-smart-switch/tl-sg105e/}
  \item Omada複数ネットワークとACL：\url{https://www.tp-link.com/jp/support/faq/3061/}
  \item Omada router-on-a-stick：\url{https://www.tp-link.com/jp/support/faq/3307/}
  \item Nginx proxy module：\url{https://nginx.org/en/docs/http/ngx_http_proxy_module.html}
  \item Nginx client\_max\_body\_size：\url{https://nginx.org/en/docs/http/ngx_http_core_module.html#client_max_body_size}
  \item Raspberry Pi UFW：\url{https://www.raspberrypi.com/documentation/security/ufw.html}
\end{itemize}
\end{sloppypar}

\newpage
\section{段階2：VLANネットワークを構築する}

\subsection{VLANとアドレス設計}

\begin{longtable}{@{}p{18mm}p{30mm}p{48mm}p{48mm}@{}}
\toprule
VLAN & 用途 & ネットワーク & 固定アドレス例\\
\midrule
20 & DMZ & \texttt{192.168.20.0/24} & P0 \texttt{192.168.20.10}\\
30 & 認証 & \texttt{192.168.30.0/24} & P1 \texttt{192.168.30.10}\\
40 & コンパイル & \texttt{192.168.40.0/24} & P2 \texttt{192.168.40.10}\\
99 & 管理 & \texttt{192.168.99.0/24} & 管理用Mac\\
\bottomrule
\end{longtable}

VLANルーターの各ゲートウェイを\texttt{192.168.X.1}とする。

\subsection{スイッチのポート設計}

\begin{longtable}{@{}p{18mm}p{38mm}p{58mm}p{20mm}@{}}
\toprule
ポート & 接続先 & VLAN設定 & PVID\\
\midrule
1 & VLANルーター & VLAN 20/30/40/99 Tagged & 99\\
2 & P0 & VLAN 20 Untagged & 20\\
3 & P1 & VLAN 30 Untagged & 30\\
4 & P2 & VLAN 40 Untagged & 40\\
5 & 管理用Mac & VLAN 99 Untagged & 99\\
\bottomrule
\end{longtable}

P0、P1、P2は通常のタグなしEthernetを使用し、VLANタグはスイッチ側で付け外しする。

\subsection{VLAN間ACL}

許可規則を拒否規則より先に置く。

\begin{longtable}{@{}p{34mm}p{38mm}p{31mm}p{39mm}@{}}
\toprule
送信元 & 宛先 & ポート & 動作\\
\midrule
管理 VLAN 99 & P0/P1/P2 & TCP 22 & 許可\\
P0 & P1 & TCP 8000 & 許可\\
P1 & P2 & TCP 8000 & 許可\\
DMZ VLAN 20 & 認証 VLAN 30 & 上記以外 & 拒否\\
DMZ VLAN 20 & コンパイル VLAN 40 & すべて & 拒否\\
認証 VLAN 30 & コンパイル VLAN 40 & 上記以外 & 拒否\\
サーバー VLAN & 管理 VLAN 99 & すべて & 拒否\\
\bottomrule
\end{longtable}

OS更新に必要なDNS、NTP、HTTPSの外向き通信は動作確認しながら許可する。二重ルーター構成では、
サーバーVLANから既存家庭内LANへの通信も拒否する。IPv6を使用する場合はIPv6にも同等の
規則を設定し、未設定の間はサーバーVLANへのグローバルIPv6配布を止める。

\section{段階2：P0を準備する}

\subsection{OSと管理ユーザー}

P0へ64-bit Raspberry Pi OS LiteまたはDebianを導入し、固定アドレス
\texttt{192.168.20.10}を設定する。初期設定後、更新する。

\begin{command}
sudo apt update
sudo apt full-upgrade
sudo reboot
\end{command}

管理用SSHはVLAN 99からだけ許可し、インターネットへ22番を転送しない。

\subsection{NginxとUFWの導入}

\begin{command}
sudo apt install nginx ufw
sudo ufw allow from 192.168.99.0/24 to any port 22 proto tcp
sudo ufw allow 443/tcp
sudo ufw default deny incoming
sudo ufw default allow outgoing
sudo ufw enable
sudo ufw status verbose
\end{command}

公開前のLAN試験で80番が必要な期間だけ、管理VLANまたはLANから限定して許可する。

\begin{command}
sudo ufw allow from 192.168.99.0/24 to any port 80 proto tcp
\end{command}

\section{段階2：NginxをP1からP0へ移す}

\subsection{設定ファイルの転送}

リポジトリで管理するP0用設定をMacからP0へ転送する。この設定はP1のVLANアドレスを
転送先とし、公開APIを3パスへ限定している。秘密鍵はこの時点では含まれない。

\begin{command}
scp Server/P0/progs/texnote-p0.nginx \
  mat@192.168.20.10:/home/mat/
\end{command}

P0上で配置する。

\begin{command}
sudo install -o root -g root -m 644 \
  /home/mat/texnote-p0.nginx \
  /etc/nginx/sites-available/texnote-p0
\end{command}

配置後、\texttt{location = /login}、\texttt{location = /me}、
\texttt{location = /compile}だけに\texttt{proxy\_pass}があり、
\texttt{location /}が404を返すことを確認する。

設定名を有効化する。

\begin{command}
sudo ln -s \
  /etc/nginx/sites-available/texnote-p0 \
  /etc/nginx/sites-enabled/texnote-p0
sudo unlink /etc/nginx/sites-enabled/default
sudo nginx -t
sudo systemctl reload nginx
\end{command}

\subsection{P1の待受アドレスをVLAN側へ変更}

P1の\texttt{tex-auth.service}を編集する。

\begin{command}
ExecStart=/home/mat/tex-auth/.auth-env/bin/uvicorn main:app --host 192.168.30.10 --port 8000
\end{command}

P1のUFWに以前追加した8000番の全面拒否を削除し、P0だけを許可する。番号は実機で確認する。

\begin{command}
sudo ufw status numbered
sudo ufw delete <8000/tcp全面拒否規則の番号>
sudo ufw allow from 192.168.20.10 to any port 8000 proto tcp
sudo ufw status numbered
\end{command}

systemd設定を反映する。

\begin{command}
sudo systemctl daemon-reload
sudo systemctl restart tex-auth.service
systemctl status tex-auth.service --no-pager -l
sudo ss -ltnp | grep ':8000'
\end{command}

\subsection{P2の接続制限}

P2はP1からのTCP 8000だけを許可する。

\begin{command}
sudo apt install ufw
sudo ufw allow from 192.168.99.0/24 to any port 22 proto tcp
sudo ufw allow from 192.168.30.10 to any port 8000 proto tcp
sudo ufw default deny incoming
sudo ufw default allow outgoing
sudo ufw enable
sudo ufw status verbose
\end{command}

P1の環境ファイルでP2の新アドレスを設定する。

\begin{command}
TEXNOTE_P2_TYPESET_URL=http://192.168.40.10:8000/v1/typeset
\end{command}

P1を再起動する。

\begin{command}
sudo systemctl restart tex-auth.service
\end{command}

\subsection{P0経由のLAN試験}

管理用Macから、P0が公開対象外のパスを拒否することを先に確認する。

\begin{command}
test "$(curl -sS -o /dev/null -w '%{http_code}' \
  http://192.168.20.10/openapi.json)" = 404
test "$(curl -sS -o /dev/null -w '%{http_code}' -X POST \
  http://192.168.20.10/compile-test)" = 404
\end{command}

続いてP0経由でログイン、\texttt{/me}、実際の画像付き\texttt{/compile}を確認する。
\texttt{/compile-test}によるP1--P2間の保守確認が必要な場合はP1へSSH接続し、
P1自身の\texttt{127.0.0.1:8000/compile-test}を使用する。P0からP2へ直接接続できず、
P0からP1のTCP 8000だけへ接続できることも確認する。

\section{段階2：HTTPSと外部公開}

\subsection{公開前提}

HTTPS証明書を取得するには、通常は所有するドメイン名、DNS設定、外部からP0へ到達できる
経路が必要である。契約回線がCGNATの場合は通常のポート転送が使えないことがあるため、
公開前に確認する。

ルーターの転送はTCP 443だけとする。

\begin{command}
Internet TCP 443 -> P0 192.168.20.10 TCP 443
\end{command}

TCP 22、80、8000を恒常的にインターネットへ転送しない。

\subsection{証明書取得}

ドメインとDNSを準備した後、P0へCertbotのNginx連携を導入する例を示す。

\begin{command}
sudo apt update
sudo apt install certbot python3-certbot-nginx
sudo certbot --nginx -d texnote.example.com
\end{command}

\texttt{texnote.example.com}は実際のドメインへ置き換える。取得後に自動更新を確認する。

\begin{command}
systemctl status certbot.timer --no-pager
sudo certbot renew --dry-run
sudo nginx -t
\end{command}

TeXNoteのP1 URLをHTTPSへ変更する。

\begin{command}
https://texnote.example.com
\end{command}

\section{P0移行後の切戻し}

P0またはVLANで問題が発生した場合、外部公開を停止してLAN内のP1 Nginxへ戻す。

\begin{enumerate}
  \item ルーターのTCP 443転送を停止する。
  \item P0 Nginxを停止する。
  \item P1 Uvicornを\texttt{127.0.0.1:8000}へ戻す。
  \item P1上のNginxを再度有効化する。
  \item TeXNoteの接続先をP1のLAN URLへ戻す。
\end{enumerate}

\begin{command}
# P0
sudo systemctl stop nginx

# P1
sudoedit /etc/systemd/system/tex-auth.service
# ExecStartのhostを127.0.0.1へ戻す
sudo systemctl daemon-reload
sudo systemctl restart tex-auth.service
sudo systemctl enable --now nginx
sudo nginx -t
\end{command}

\section{完了チェックリスト}

\subsection{段階1：P1上のリバースプロキシ}

\begin{itemize}
  \item[$\square$] Nginxがenabledかつactiveである。
  \item[$\square$] Nginxの要求上限が32 MiBである。
  \item[$\square$] P1 Uvicornが\texttt{127.0.0.1:8000}だけで待ち受けている。
  \item[$\square$] LANからTCP 8000へ直接接続できない。
  \item[$\square$] Nginx経由でlogin、me、compileが成功する。
  \item[$\square$] UFW有効化後も新しいSSH接続が成功する。
  \item[$\square$] P1再起動後もNginxとP1が自動復旧する。
\end{itemize}

\subsection{段階2：P0とVLAN}

\begin{itemize}
  \item[$\square$] P0、P1、P2、管理用端末が別VLANへ分離された。
  \item[$\square$] P0からP1のTCP 8000だけが許可された。
  \item[$\square$] P1からP2のTCP 8000だけが許可された。
  \item[$\square$] P0からP2へ直接接続できない。
  \item[$\square$] インターネットへ公開されるのはP0のTCP 443だけである。
  \item[$\square$] SSHは管理VLANからだけ許可される。
  \item[$\square$] IPv6にも同等の制限がある、またはサーバーVLANで無効化されている。
  \item[$\square$] HTTPS証明書の自動更新試験が成功した。
  \item[$\square$] P0、P1、P2再起動後に全経路が自動復旧する。
  \item[$\square$] 段階1へ戻す切戻し手順を確認した。
\end{itemize}

\section{まとめ}

最初にP1上でNginxを構築すると、TeXNote、P1 FastAPI、P2 FastAPIの処理を変えずに
リバースプロキシ経由の動作を確認できる。次にVLAN機器とP0を追加し、Nginxの
\texttt{proxy\_pass}をlocalhostからP1の内部アドレスへ変更することで、公開境界だけを
P0へ移せる。安全性を決めるのは機器を分けることだけではなく、P0からP1、P1からP2という
必要最小限の通信だけをACLと各ホストのUFWで許可することである。

\end{document}
